\documentclass[11pt]{article}
\usepackage[a4paper,margin=0.86in]{geometry}
\usepackage{amsmath,amssymb,amsthm,mathtools}
\usepackage{booktabs,array}
\usepackage{graphicx}
\usepackage{xcolor}
\usepackage{xurl}
\usepackage[hidelinks]{hyperref}
\hypersetup{
  pdftitle={The Action-Memory Diamond: Exact Resource Regions for Passive Matrix Networks},
  pdfauthor={Lluis Eriksson},
  pdfsubject={Exact attainable regions for arc action and McMillan degree}
}
\usepackage[nameinlink,capitalise]{cleveref}
\usepackage{microtype}

\newtheorem{theorem}{Theorem}[section]
\newtheorem{proposition}[theorem]{Proposition}
\newtheorem{lemma}[theorem]{Lemma}
\newtheorem{corollary}[theorem]{Corollary}
\theoremstyle{definition}
\newtheorem{definition}[theorem]{Definition}
\newtheorem{example}[theorem]{Example}
\theoremstyle{remark}
\newtheorem{remark}[theorem]{Remark}
\newtheorem{limitation}[theorem]{Limitation}

\newcommand{\C}{\mathbb C}
\newcommand{\D}{\mathbb D}
\newcommand{\T}{\mathbb T}
\newcommand{\tr}{\operatorname{tr}}
\newcommand{\rank}{\operatorname{rank}}
\newcommand{\ran}{\operatorname{ran}}
\newcommand{\diag}{\operatorname{diag}}
\newcommand{\dd}{\mathrm d}
\newcommand{\e}{\mathrm e}
\newcommand{\mcdeg}{\deg_{\mathrm{McM}}}
\newcommand{\Gr}{\operatorname{Gr}}

\title{The Action--Memory Diamond:\\
Exact Resource Regions for Passive Matrix Networks}
\author{Lluis Eriksson}
\date{10 August 2026}

\begin{document}
\maketitle

\begin{abstract}
A passive lossless network pays two different costs when it rotates a signal
subspace over a frequency arc: a continuous Wigner--Smith action and an
integer McMillan memory.  Known Grassmannian length bounds and rational-inner
degree identities constrain these resources separately.  For the stated
finite rational-inner class, we determine their joint attainable region
exactly.  Let
two rank-$k$ subspaces have principal angles $\beta_j$, put
$B=\sum_j\beta_j$ and let $r$ be the number of nonzero angles.  For a square
finite rational inner transfer matrix of degree at most $d$, define the arc
action
\[
 A=\int_I\tr Q(\theta)\,\dd\theta,
 \qquad Q=-\mathrm iS^*\partial_\theta S\succeq0.
\]
If $B>0$, the exact feasible set is
\[
 \boxed{\ d\ge r,\qquad 2B\le A\le2\pi d-2B\ }.
\]
Both faces are attained.  Equivalently,
\[
 d_{\min}(A)=\max\!\left\{r,
 \left\lceil\frac{A+2B}{2\pi}\right\rceil\right\}.
\]
The upper face is a return cost: the complementary arc must rotate the
subspace back, while the full-circle trace action is exactly $2\pi$ times
the degree.  We prove sufficiency by an explicit compiler of normalized
rank-one Blaschke--Potapov gates; it realizes every interior point and both
faces, including rank-deficient and orthogonal cases.  When the endpoint
subspaces coincide, the region changes discontinuously to
$A\in[0,2\pi d)$, with the upper endpoint open.

We derive fail-closed noisy degree certificates and give a multi-frequency
warning using classical boundary Nevanlinna--Pick theory: three pairwise
degree-one routing tasks can require degree two jointly.  For orthogonal-line
data the obstruction is a frame-independent cycle phase, exhibited by an exact positive semidefinite
Pick completion with spectrum $(2,1,0)$.  Deterministic code compiles random
points of the diamond, constructs the degree-two colligation, and is replayed
by an independent verifier.
\end{abstract}

\noindent\textbf{Keywords:} rational inner matrix; Wigner--Smith delay;
McMillan degree; Blaschke--Potapov product; Grassmannian; passive network;
boundary Nevanlinna--Pick interpolation; resource theory.

\section{One routing event, two coupled resources}

For a lossless causal finite-dimensional network, the boundary transfer
matrix $S(\e^{\mathrm i\theta})$ is unitary and its proper-delay matrix
\begin{equation}\label{eq:delay}
 Q(\theta)=-\mathrm iS(\e^{\mathrm i\theta})^*
 \frac{\dd}{\dd\theta}S(\e^{\mathrm i\theta})
\end{equation}
is positive semidefinite.  The integral of $\tr Q$ over an arc measures a
continuous dwell-time or phase-action resource.  The McMillan degree counts
the dimension of a minimal conservative realization and is integer valued
\cite{McMillan1952,GoughZhang2015}.

These resources cannot be chosen independently.  A subspace rotation costs
action on the observed arc.  Less obviously, the same network must close its
unitary loop on the complementary arc.  The trace action available to pay
both journeys is quantized by the determinant winding.  This observation
turns two one-sided lower bounds into the closed diamond in
\cref{fig:diamond}.

Grassmann metrics and optimal unitary/Grassmann paths are classical
\cite[Thm.~1]{QiuZhangLi2005}
\cite[Thms.~16 and 20]{AntezanaLarotondaVarela2014}; see also the recent
projector formulation \cite[Sec.~IV.4]{Huang2026}.  Rank-one
Blaschke--Potapov factorization and boundary all-pass interpolation are
likewise classical \cite[Thm.~3.1]{AlpayJorgensenLewkowicz2016}
\cite[Thms.~4.7, 5.1 and 8.2]{BolotnikovDym2006}
\cite[Thm.~1 and Sec.~4.4]{BharathEtAl2023}.  General rational interpolation
with stability and minimum-degree constraints is also classical
\cite{BallKim1994}.  Within the sources compared in \cref{tab:precedent}, we
did not find an exact joint attainable region for arc-integrated trace delay
and McMillan degree under a prescribed subspace transfer.  That conjunction,
including both converses and all boundary cases, is the contribution here;
this is a scoped literature statement, not a claim that the surrounding
interpolation theories are new.

\begin{figure}[t]
 \centering
 \includegraphics[width=0.88\linewidth]{figures/action_memory_diamond.pdf}
 \caption{Left: the exact action--memory region for fixed endpoint geometry.
 The lower face is the forward Grassmann cost and the upper face reserves the
 same cost for the complementary return.  Right: a three-node
 Grassmann--Pick completion has rank two although every pair has rank one;
 its cycle product has phase $\pi/2$.}
 \label{fig:diamond}
\end{figure}

\paragraph{What is and is not claimed.}
The theorem concerns square finite rational inner matrices.  It does not
cover loss, gain, singular inner factors, nonrational continuous-time
propagation delays or a
non-Markovian bath.  The quantity $A$ is trace delay, not thermodynamic work.
The multi-node rank formula used later is a gauge-quotient corollary of
classical tangential boundary interpolation, not a new Pick theorem.

\begin{table}[t]
\centering
\small
\begin{tabular}{@{}>{\raggedright\arraybackslash}p{0.22\linewidth}
>{\raggedright\arraybackslash}p{0.31\linewidth}
>{\raggedright\arraybackslash}p{0.37\linewidth}@{}}
\toprule
Precedent & Established ingredient & Not addressed there \\
\midrule
Qiu--Zhang--Li \cite{QiuZhangLi2005}
& Symmetric-gauge Grassmann metrics
& Positive delay, rational degree and complementary-arc budget \\
Antezana--Larotonda--Varela \cite{AntezanaLarotondaVarela2014}
& Optimal unitary and Grassmann paths
& Inner analyticity, winding quantization and converse synthesis \\
Potapov; Alpay et al.\ \cite{Potapov1955,AlpayJorgensenLewkowicz2016}
& Rank-one factorization and degree
& Exact action region under endpoint subspace data \\
Bolotnikov--Dym \cite{BolotnikovDym2006}
& Matrix tangential boundary interpolation
& Arc-integrated action and forward/return geometry \\
Ball--Kim \cite{BallKim1994}
& Stability and McMillan degree for rational interpolants
& Inner subspace endpoints and a prescribed arc-action region \\
Bharath et al.\ \cite{BharathEtAl2023}
& All-pass design with nodewise group-delay optimization
& Closed action--degree diamond and its two boundary compilers \\
\bottomrule
\end{tabular}
\caption{Closest precedents and scope distinction.  Each cited work supplies
the listed ingredient; we are not aware of one containing the present joint
attainable-region theorem.}
\label{tab:precedent}
\end{table}

\begin{table}[t]
\centering
\small
\begin{tabular}{@{}>{\raggedright\arraybackslash}p{0.23\linewidth}
>{\raggedright\arraybackslash}p{0.33\linewidth}
>{\raggedright\arraybackslash}p{0.34\linewidth}@{}}
\toprule
Question & Preceding paper \cite{ErikssonDelayRegions2026} & Present paper \\
\midrule
Continuous two-node data
& Exact modal region $2\beta\prec_w q$ for absolutely continuous positive-generator paths
& No new modal-profile claim; scalar trace action in the finite rational-inner class \\
Rational two-node memory
& Minimum degree $r$ and a construction on the lower face $A=2B$
& Complete iff region in $(A,d)$, including the upper return face, all interior points and the discontinuous $B=0$ case \\
Multi-node information
& Pick--Wigner--Smith Gram/rank theorem for specified unitary values and local derivatives
& Gauge-quotiented subspace corollary and an exact three-line cyclic obstruction; no new general Pick theorem \\
Robust conclusion
& Separate lower certificates for modal action and hidden-state rank
& Joint feasibility/falsification for a proposed $(A,d)$ budget and a cycle-product null control \\
\bottomrule
\end{tabular}
\caption{Theorem-level non-overlap with the immediately preceding paper.
The earlier modal theorem is finer in spectral profile; the present theorem
is stronger only in its exact coupling of total arc action to integer degree.}
\label{tab:companion}
\end{table}

\section{Setup}

Let $I=[\theta_0,\theta_1]$ be a proper positively oriented arc, let
$d\in\mathbb N_0$, and suppose
$0<\theta_1-\theta_0<2\pi$, and write
$\zeta_j=\e^{\mathrm i\theta_j}$.  Let $S$ be an $N$-by-$N$ rational inner
matrix regular on $\T$.  Its McMillan degree is denoted $n=\mcdeg S$.

Fix $\mathcal X\in\Gr(k,N)$ and set
\[
 \mathcal Y_0=S(\zeta_0)\mathcal X,
 \qquad
 \mathcal Y_1=S(\zeta_1)\mathcal X.
\]
Let
\[
 \frac{\pi}{2}\ge\beta_1\ge\cdots\ge\beta_k\ge0
\]
be their canonical angles and define
\begin{equation}\label{eq:Br}
 B=\sum_{j=1}^k\beta_j,
 \qquad
 r=\#\{j:\beta_j>0\}.
\end{equation}
The arc action is
\begin{equation}\label{eq:action}
 A_I(S)=\int_{\theta_0}^{\theta_1}\tr Q(\theta)\,\dd\theta.
\end{equation}

For a finite rational inner matrix, $Q(\theta)\succeq0$ on $\T$ and
\begin{equation}\label{eq:winding}
 \int_0^{2\pi}\tr Q(\theta)\,\dd\theta
 =2\pi\,\operatorname{wind}(\det S)
 =2\pi\mcdeg S.
\end{equation}
The last equality follows directly from a rank-one
Blaschke--Potapov factorization \cite{Potapov1955,
AlpayJorgensenLewkowicz2016}.

\section{The exact action--memory diamond}

\begin{theorem}[Exact action--memory diamond]\label{thm:diamond}
Let $I$ be a proper arc and let $\mathcal X,\mathcal Y\in\Gr(k,N)$ have
principal-angle data $(B,r)$ from \eqref{eq:Br}.

If $B>0$, there exists an $N$-by-$N$ finite rational inner matrix $S$,
regular on $\T$, such that
\[
 S(\zeta_0)\mathcal X=\mathcal X,
 \qquad S(\zeta_1)\mathcal X=\mathcal Y,
 \qquad \mcdeg S\le d,
 \qquad A_I(S)=A,
\]
if and only if
\begin{equation}\label{eq:diamond}
 \boxed{d\ge r,\qquad 2B\le A\le2\pi d-2B.}
\end{equation}
Constants on the left allow an arbitrary prescribed initial image in place
of $\mathcal X$.

If $B=0$, the exact region is
\begin{equation}\label{eq:trivialregion}
 A=0\quad\text{or}\quad d\ge1\text{ and }0<A<2\pi d.
\end{equation}
Thus for $d\ge1$ it is $[0,2\pi d)$; for $d=0$ it is the singleton
$\{0\}$.
\end{theorem}

\begin{corollary}[Exact least memory at prescribed action]
\label{cor:dmin}
For $B>0$ and $A\ge2B$,
\begin{equation}\label{eq:dmin}
 d_{\min}(A)=
 \max\left\{r,
 \left\lceil\frac{A+2B}{2\pi}\right\rceil\right\}.
\end{equation}
For $B=0$, $d_{\min}(0)=0$ and
\[
 d_{\min}(A)=\left\lfloor\frac{A}{2\pi}\right\rfloor+1,
 \qquad A>0.
\]
\end{corollary}

Three features are worth isolating.  First, the two faces in
\eqref{eq:diamond} are inclusive even when some angles are orthogonal.
Second, an arbitrarily small but positive angle forces one unit of memory,
whereas at $B=0$ the degree can drop to zero.  Third, the open upper endpoint
in \eqref{eq:trivialregion} is real: every nonconstant finite inner function
has strictly positive trace delay on every nonempty complementary arc.

\section{Necessity: the return arc spends the missing action}

We give a self-contained trace version of the geometric speed limit.

\begin{lemma}[Positive-generator subspace speed]\label{lem:speed}
Let $U(t)$ be a differentiable unitary path with
$-\mathrm iU^*\dot U=H(t)\succeq0$, and let
$\mathcal X(t)=U(t)\mathcal X(0)$.  If the endpoint principal angles are
$\beta_j$, then
\begin{equation}\label{eq:trace-speed}
 2\sum_j\beta_j\le\int\tr H(t)\,\dd t.
\end{equation}
\end{lemma}

\begin{proof}
Let $P$ be the moving orthogonal projector and write $H$ in blocks relative
to $P\oplus(I-P)$ as
$\left(\begin{smallmatrix}H_{11}&K\\K^*&H_{22}\end{smallmatrix}\right)$.
The projector velocity has trace norm
$\lVert\dot P\rVert_1=2\lVert K\rVert_*$.  Positivity of the block matrix
implies
\[
 2\lVert K\rVert_*\le\tr H_{11}+\tr H_{22}=\tr H;
\]
for example, test positivity against the polar partial isometry of $K$.
The trace-norm Grassmann length dominates twice the sum of endpoint principal
angles \cite{QiuZhangLi2005,AntezanaLarotondaVarela2014}.  Integration gives
\eqref{eq:trace-speed}.
\end{proof}

Applying \cref{lem:speed} on $I$ gives $A\ge2B$.  On the complementary arc,
the same periodic transfer matrix returns $\mathcal Y$ to $\mathcal X$.
Its action is $2\pi n-A$ by \eqref{eq:winding}, so
\begin{equation}\label{eq:return}
 2\pi n-A\ge2B.
\end{equation}
Since $n\le d$, this yields the upper face in \eqref{eq:diamond}.

It remains to see why $n\ge r$.  Factor $S$, up to a constant unitary, into
$n$ rank-one Blaschke--Potapov factors $F_j$.  Put
$P_j(z)=F_1(z)\cdots F_j(z)$ and
$R_j=P_j(\zeta_1)P_j(\zeta_0)^*$.  Then
\[
 R_j=\left[P_{j-1}(\zeta_1)F_j(\zeta_1)F_j(\zeta_0)^*
 P_{j-1}(\zeta_1)^*\right]R_{j-1}.
\]
The bracket is a unitary gate $I+(\lambda-1)vv^*$ after conjugating its
rank-one projector.  Such a gate can
lower the dimension of the intersection of two $k$-planes by at most one.
After $n$ gates,
\[
 \dim(\mathcal X\cap\mathcal Y)\ge k-n.
\]
Because $r=k-\dim(\mathcal X\cap\mathcal Y)$, one has $n\ge r$.

When $B=0$, \eqref{eq:winding} only gives $A\le2\pi n$.  Equality cannot
hold for $n>0$: the complementary arc is nonempty and each scalar Poisson
kernel in a Potapov factorization is strictly positive there.  Hence
$A<2\pi n\le2\pi d$.  The constant case gives $A=0$.

\section{Two exact one-gate lemmas}

The converse rests on two elementary facts.  We state them explicitly to
fix phases and endpoints.

\begin{lemma}[Any proper-arc phase]\label{lem:phase}
For a proper arc $I$ and every $\delta\in(0,2\pi)$, there is a scalar
degree-one Blaschke factor $b$ such that
\[
 b(\zeta_0)=1,
 \qquad b(\zeta_1)=\e^{\mathrm i\delta},
 \qquad \int_I-\mathrm i\overline b\,\partial_\theta b\,\dd\theta=\delta.
\]
\end{lemma}

\begin{proof}
Rotate the arc to $[-\gamma,\gamma]$, where
$\gamma=(\theta_1-\theta_0)/2\in(0,\pi)$.  For
$b_\rho(z)=(z-\rho)/(1-\rho z)$, $-1<\rho<1$, its phase advance is
\begin{equation}\label{eq:phaseformula}
 \delta_\rho=4\arctan\left(
 \frac{1+\rho}{1-\rho}\tan\frac{\gamma}{2}\right).
\end{equation}
It increases continuously from $0$ to $2\pi$.  More explicitly, the desired
zero is
\begin{equation}\label{eq:rhofromdelta}
 \rho=
 \frac{\tan(\delta/4)-\tan(\gamma/2)}
      {\tan(\delta/4)+\tan(\gamma/2)}.
\end{equation}
Multiplication by a unimodular constant normalizes the value at the first
endpoint.  The integral equals the unwrapped phase advance.
\end{proof}

\begin{lemma}[Exact rank-one line gate]\label{lem:gate}
Let two lines have angle $\alpha\in(0,\pi/2]$.  A rank-one unitary
\[
 W=I+(\e^{\mathrm i\delta}-1)vv^*
\]
maps the first line to the second, after choosing $v$, if and only if
\begin{equation}\label{eq:gateinterval}
 2\alpha\le\delta\le2\pi-2\alpha.
\end{equation}
\end{lemma}

\begin{proof}
In the plane of the two lines, write
$v=\cos\eta\,x+\e^{\mathrm i\varphi}\sin\eta\,x_\perp$.  Then
\begin{equation}\label{eq:gateoverlap}
 |\langle x,Wx\rangle|^2
 =1-\sin^2(\delta/2)\sin^2(2\eta).
\end{equation}
Thus $|\langle x,Wx\rangle|=\cos\alpha$ is possible exactly when
$\sin\alpha\le|\sin(\delta/2)|$, which is \eqref{eq:gateinterval}.
For a constructive choice, set
\[
 t=\cos^2\eta=
 \frac{1-\sqrt{1-\sin^2\alpha/\sin^2(\delta/2)}}{2}
\]
and choose $\varphi$ to match the phase of the prescribed second line.
This also covers both equalities in \eqref{eq:gateinterval}.
\end{proof}

\section{Sufficiency: compiling every point}

Assume first that $B>0$.  Choose mutually orthogonal principal planes with
vectors $x_j\in\mathcal X$ and $x_j^\perp\in\mathcal X^\perp$ such that the
target principal line is
\[
 y_j=\cos\beta_j\,x_j+\sin\beta_j\,x_j^\perp.
\]
Only the $r$ positive angles need gates.

Fix an integer $n$ with $r\le n\le d$ and
\begin{equation}\label{eq:fixedn}
 2B\le A\le2\pi n-2B.
\end{equation}
Choose positive integers $q_j$ summing to $n$ and subdivide each principal
geodesic into positive increments
\[
 \beta_j=\alpha_{j1}+\cdots+\alpha_{jq_j}.
\]
By \cref{lem:gate}, the phase assigned to step $(j,\ell)$ may be any value
in
\[
 [2\alpha_{j\ell},\,2\pi-2\alpha_{j\ell}].
\]
The Minkowski sum of these $n$ intervals is exactly the interval in
\eqref{eq:fixedn}.  Hence phases $\delta_{j\ell}$ can be chosen to sum to
$A$.

Use \cref{lem:phase} to produce normalized scalar factors
$b_{j\ell}$ with these arc phases.  At each intermediate principal line,
use \cref{lem:gate} to choose a rank-one projector $v_{j\ell}v_{j\ell}^*$,
and define
\[
 B_{j\ell}(z)=I+(b_{j\ell}(z)-1)v_{j\ell}v_{j\ell}^*.
\]
Multiply the factors in the physical order of the intermediate lines.
Every factor is the identity at $\zeta_0$, while at $\zeta_1$ the gates
carry each $x_j$ along its principal plane to $y_j$.  The untouched
intersection is fixed.  Thus the product maps $\mathcal X$ to $\mathcal Y$.

The product is rational inner.  Its delay is a sum of unitary conjugates of
the positive rank-one scalar delays, so it is positive.  Its determinant is
the product of the $n$ scalar factors; consequently its degree is exactly
$n$, its full action is $2\pi n$, and its arc action is
$\sum_{j,\ell}\delta_{j\ell}=A$.  This proves sufficiency for a fixed $n$.
The union over $r\le n\le d$ is the entire diamond
\eqref{eq:diamond}: the intervals are nested apart from their already shared
lower endpoint, which is attained at degree $r$.

If $B=0$ and $A>0$, choose
$n=\lfloor A/(2\pi)\rfloor+1\le d$.  Split $A$ into $n$ numbers in
$(0,2\pi)$ and place all corresponding normalized scalar factors on a fixed
line inside $\mathcal X$.  The endpoint unitary preserves $\mathcal X$ and
the action is $A$.  The constant transfer handles $A=0$.  This completes the
proof of \cref{thm:diamond} and \cref{cor:dmin}.

\begin{remark}[An explicit compiler]
The proof is algorithmic.  Equal subdivisions
$\alpha_{j\ell}=\beta_j/q_j$ suffice.  Start every phase at its lower value
$2\alpha_{j\ell}$ and distribute the surplus $A-2B$ across capacities
$2\pi-4\alpha_{j\ell}$.  Equations \eqref{eq:rhofromdelta} and
\eqref{eq:gateoverlap} then return the Blaschke zero and projector without
optimization.
\end{remark}

\section{Sharp examples and discontinuities}

\paragraph{A collapsed diamond.}
For one orthogonal pair, $\beta=\pi/2$, $r=d=1$.  The lower and upper faces
coincide, so $A=\pi$ is forced.  More generally, $r=d$ mutually orthogonal
principal lines force $A=\pi r$.  Half the winding is spent on the forward
arc and half on the return.

\paragraph{Extra memory widens by $2\pi$.}
For $\beta=(\pi/3,\pi/6,0)$, $B=\pi/2$ and $r=2$.  Degree two gives
$A\in[\pi,3\pi]$; degree three gives $A\in[\pi,5\pi]$.  The extra state does
not reduce the geometric minimum, but it adds one full turn to the available
upper action.

\paragraph{Topological jump at zero motion.}
For every $\beta>0$, however small, $d_{\min}(2\beta)=1$.  At $\beta=0$ and
$A=0$, it is zero.  Continuous action can vanish along a sequence while the
integer memory remains fixed.  Conversely, when $B=0$ and $A$ approaches
$2\pi d$, the required poles concentrate near the boundary but the endpoint
subspace still does not move.

\section{Fail-closed inference from noisy data}

Suppose a projector experiment supplies certified lower bounds
$\underline\beta_j\le\beta_j$, and an action measurement supplies
$A\in[\underline A,\overline A]$.  Put
\[
 \underline B=\sum_j\underline\beta_j,
 \qquad
 \underline r=\#\{j:\underline\beta_j>0\}.
\]

\begin{corollary}[Robust memory certificate]\label{cor:robust}
Every compatible rational inner model satisfies
\begin{equation}\label{eq:robustdegree}
 \mcdeg S\ge
 \max\left\{\underline r,
 \left\lceil\frac{\underline A+2\underline B}{2\pi}\right\rceil
 \right\},
\end{equation}
whenever $\underline B>0$.  A claimed degree budget $d$ is falsified if
\[
 \overline A<2\underline B
 \quad\text{or}\quad
 \underline A>2\pi d-2\underline B.
\]
\end{corollary}

The certificate is deliberately one-sided.  It never substitutes point
estimates for uncertainty sets, and therefore cannot overstate the required
degree when the stated error bounds hold.  Standard perturbation bounds for
projectors turn operator-norm error bars into lower principal-angle bounds;
the preceding paper gives one such implementation
\cite{ErikssonDelayRegions2026}.

\section{Beyond two nodes: a cyclic memory invisible pairwise}

The diamond completely solves one arc with two endpoint subspaces.  It does
not say that pairwise costs determine a multi-frequency interpolation.  The
right classical object is a boundary Pick completion.

Fix an input frame $X\in\C^{N\times k}$, boundary nodes $\zeta_i$, target
frames $Y_i$, and gauges $V_i\in U(k)$.  Put $W_i=Y_iV_i$ and prescribe
off-diagonal blocks
\begin{equation}\label{eq:gp}
 P_{ij}=\frac{I_k-W_i^*W_j}{1-\overline\zeta_i\zeta_j},
 \qquad i\ne j,
\end{equation}
while leaving Hermitian diagonal blocks free.  Tangential boundary
Nevanlinna--Pick theory implies that the least rational-inner degree equals
the least rank of a positive semidefinite completion, minimized also over
the gauges \cite[Thms.~4.7, 5.1 and 8.2]{BolotnikovDym2006}.  Equivalently,
this follows from a lurking-isometry construction and the compressed
de Branges--Rovnyak kernel.  We use this classical corollary only to expose a
strict global effect.

\begin{proposition}[Three-node cyclic memory]\label{prop:cycle}
Let $N=3$, $k=1$, $\omega=\e^{2\pi\mathrm i/3}$, and take nodes
$1,\omega,\omega^2$ with three mutually orthogonal target lines.  Every
two-node subproblem has minimum degree one, but the joint problem has minimum
degree two.
\end{proposition}

\begin{proof}
All off-diagonal target overlaps vanish, independently of scalar frame
phases.  With
\[
 a=\frac{1}{1-\omega}=\frac12+\frac{\mathrm i}{2\sqrt3},
\]
the diagonal choice $P_{00}=P_{11}=P_{22}=1$ gives
\begin{equation}\label{eq:cycleP}
 P=\begin{pmatrix}
 1&a&\overline a\\
 \overline a&1&a\\
 a&\overline a&1
 \end{pmatrix},
 \qquad \operatorname{spec}P=(2,1,0).
\end{equation}
Thus a degree-two interpolant exists.  If a positive semidefinite rank-one
completion existed, $P_{ij}=\overline f_i f_j$ would force
$P_{01}P_{12}P_{20}=|f_0f_1f_2|^2$ to be nonnegative real.  Instead,
\begin{equation}\label{eq:cyclephase}
 P_{01}P_{12}P_{20}=a^3=\frac{\mathrm i}{3\sqrt3}.
\end{equation}
Hence rank one is impossible.  For any pair, choosing both diagonal entries
as $|P_{ij}|$ gives a positive semidefinite rank-one completion.
\end{proof}

For the present orthogonal-line data the cross blocks, and hence their cycle
product, are independent of all frame gauges because every off-diagonal
target overlap vanishes.  The phase is a holonomy-like cycle invariant and,
in this example, records information not visible in pairwise canonical
angles.  The underlying Pick theorem and its rank formulation are classical;
here the example demonstrates the strict separation
\[
 d_{\min}^{(3\ \mathrm{nodes})}=2>
 \max_{i\ne j}d_{\min}^{(i,j)}=1.
\]

\paragraph{A robust cycle witness.}
For these orthogonal data, or for fixed phase-calibrated tangential frames,
let three measured nonzero cross blocks be scalars $\widehat p_1,
\widehat p_2,\widehat p_3$, each within $\varepsilon$ of its exact value,
and let $M=\max_j|\widehat p_j|$.  Telescoping the product gives
\[
 \left|\prod_jp_j-\prod_j\widehat p_j\right|
 \le3\varepsilon(M+\varepsilon)^2.
\]
Therefore
\[
 \left|\operatorname{Im}\prod_j\widehat p_j\right|
 >3\varepsilon(M+\varepsilon)^2
\]
certifies rank at least two without selecting diagonal blocks or solving a
rank minimization.

\section{Reproducible adversarial audit}

The proofs are analytic; computation tests the formulas at their most likely
failure interfaces.  The producer and verifier are separate programs and the
verifier imports no producer functions.  The numerical compiler never rounds
a positive angle into the $B=0$ branch; inputs below $10^{-10}$ radians are
reported as numerically indeterminate.  The frozen campaign includes:
\begin{itemize}
 \item $128$ randomized feasible diamonds with ranks one through four,
 prescribed actions, analytic derivatives and adaptive quadrature, plus a
 separate noncommuting two-gate order-reversal control;
 \item exact lower and upper faces, orthogonal collapsed diamonds, zero
 angles, a degree obstruction and the open $B=0$ endpoint;
 \item $512$ independent scalar phase inversions and $2{,}048$ independent
 rank-one gate identities;
 \item the exact Pick spectrum $(2,1,0)$, cycle product
 $\mathrm i/(3\sqrt3)$, a synthesized degree-two unitary colligation and
 direct boundary interpolation checks;
 \item $1{,}024$ noisy trials around the rank-two cycle and $1{,}024$
 rank-one null controls for the fail-closed product bound.
\end{itemize}

For the flagship point
\[
 \beta=(0.92,0.47,0.18),\quad B=1.57,\quad A=15.25,
\]
the compiler uses degree three.  The observed maximum principal-angle error
is below $3\times10^{-16}$, the arc-action error is below
$2\times10^{-14}$, and the full-circle action differs from $6\pi$ by less
than $2\times10^{-14}$.  Across the randomized campaign the frozen maximum
arc-action error is below $10^{-10}$ and no positivity, endpoint or winding
test fails.  These figures are diagnostics, not evidence replacing the
proof.

\section{Physical reading and falsifiers}

The lower face $A=2B$ is a direct routing cost.  The upper face is equally
physical: a finite passive network has a quantized total phase budget, and
the unobserved frequencies must retain enough action to close the subspace
loop.  Adding one internal state enlarges the allowed upper action by exactly
$2\pi$.  The theorem therefore predicts a staircase in minimum realization
order when arc dwell time is increased at fixed endpoint geometry.

The result is falsifiable in several ways.  A finite rational inner network
inside the stated class with $A<2B$, $A>2\pi d-2B$, or $d<r$ would refute
necessity.  A feasible point for which the explicit gate compiler fails
would refute sufficiency.  Lossy, active or nonrational devices are not
counterexamples because they leave the model class.  Repeated boundary nodes
would require confluent Pick data in the multi-node section and are not
covered by \eqref{eq:gp}.

\section{Conclusion}

Continuous delay and discrete memory form one exact resource region.  For a
nontrivial subspace transfer, the forward and return arcs each pay the same
Grassmann trace cost $2B$, while the determinant winding supplies exactly
$2\pi d$.  Nothing is lost between these necessary inequalities: normalized
rank-one Blaschke--Potapov gates compile every point of the diamond.  The
least memory at prescribed action follows in closed form.

The three-node example marks the boundary of the two-point story.  Pairwise
angles see support motion; a cyclic Pick phase can demand additional hidden
state.  This suggests a next problem with genuine content: characterize
which cycle holonomies raise the minimum positive-completion rank for
matrix-valued subspace data, rather than merely restating general boundary
interpolation.

\section*{Data and code availability}

Source, deterministic certificates, the independent verifier and the exact
PDF are archived at
\href{https://github.com/lluiseriksson/finite-sample-spectral-certificates/releases/tag/v1.9-action-memory-diamond-final}{the immutable GitHub release}.  The principal entry points are
\path{research/exact_action_memory_certificate.py} and
\path{verification/verify_exact_action_memory.py}.  The certificate
contains deterministic seeds, tolerances and every reported diagnostic.

\appendix

\section{Proof details omitted from the main line}

\subsection{The positive-block trace inequality}

For completeness, let
\[
 H=\begin{pmatrix}A&K\\K^*&C\end{pmatrix}\succeq0
\]
on finite-dimensional spaces.  If $K=V|K|$ is its polar decomposition, the
standard contraction factorization gives $K=A^{1/2}TC^{1/2}$ with
$\lVert T\rVert\le1$.  Equivalently, apply
$2\operatorname{Re}\tr(V^*K)\le\tr A+\tr C$ after completing the square.
Since $\operatorname{Re}\tr(V^*K)=\tr|K|$, one obtains
\begin{equation}\label{eq:block-detail}
 2\lVert K\rVert_*\le\tr A+\tr C.
\end{equation}

For the projector path $P(t)=U(t)P_0U(t)^*$ with
$\dot U=\mathrm iUH$, direct differentiation gives
\[
 \dot P=\mathrm iU[H,P_0]U^*.
\]
Relative to $P_0\oplus(I-P_0)$, the commutator has off-diagonal blocks
$-K$ and $K^*$, hence $\lVert\dot P\rVert_1=2\lVert K\rVert_*$.  The
trace-norm length of a Grassmann path is at least
$2\sum_j\beta_j$ \cite[Thm.~1]{QiuZhangLi2005}
\cite[Thm.~20]{AntezanaLarotondaVarela2014}.  Combining this fact with
\eqref{eq:block-detail} proves
\cref{lem:speed} without a smoothness assumption beyond absolute
continuity.

\subsection{Why one Potapov factor changes at most one support direction}

Let $R=I+(\lambda-1)vv^*$ and let $\mathcal M$ be a $k$-plane.  Every vector
in $\mathcal M\cap v^\perp$ is fixed, so
\[
 \dim(\mathcal M\cap R\mathcal M)\ge k-1.
\]
For a chain $\mathcal M_j=R_j\mathcal M_{j-1}$, the codimension inside the
initial plane can increase by at most one per step.  An induction based on
\[
 \dim(E\cap G)\ge\dim(E\cap F)+\dim(F\cap G)-\dim F
\]
gives $\dim(\mathcal M_0\cap\mathcal M_n)\ge k-n$.  Thus the number of
positive canonical angles is at most the number of rank-one factors.  This
also shows why a small angle does not mean a fraction of a state.

\subsection{Endpoint regularity and strict positivity off the observed arc}

A scalar Blaschke factor with zero $a\in\D$ has boundary phase derivative
\[
 p_a(\theta)=\frac{1-|a|^2}{|\e^{\mathrm i\theta}-a|^2}>0.
\]
For a rank-one Potapov factor its delay is $p_a(\theta)vv^*$, and the delay
of a product is a sum of unitary conjugates of these terms.  Therefore a
nonconstant finite product has $\tr Q(\theta)>0$ at every boundary point.
All poles are outside $\overline\D$, so the constructed transfer is regular
at both endpoints and on the entire circle.  This proves both the open face
in \eqref{eq:trivialregion} and the absence of hidden endpoint singularities
in the compiler.

\section{A gauge-quotient boundary Pick corollary}

We record the exact classical result used in \cref{prop:cycle} so that the
degree conclusion can be audited without an appeal to engineering
terminology.

\begin{proposition}[Boundary Pick rank after quotienting frames]
\label{prop:classical-pick}
Let $\zeta_1,\ldots,\zeta_L\in\T$ be distinct, fix an isometric input frame
$X\in\C^{N\times k}$, and let $Y_i$ be isometric frames for prescribed
target $k$-planes.  For gauges $V_i\in U(k)$ put $W_i=Y_iV_i$ and form the
off-diagonal blocks \eqref{eq:gp}.  Then the minimum McMillan degree of a
square rational inner $S$, regular at the nodes and satisfying
$S(\zeta_i)\ran X=\ran Y_i$, equals
\begin{equation}\label{eq:pick-rank-min}
 \min\left\{\rank P:
 P\succeq0,\ P_{ij}\text{ obeys \eqref{eq:gp} for }i\ne j,
 \ V_i\in U(k)\right\},
\end{equation}
where the Hermitian diagonal blocks are free.
\end{proposition}

This is a quotient formulation of matrix tangential boundary interpolation
\cite{BolotnikovDym2006}; the short proof is included to expose exactly where
the rank enters.  Feasible completions exist after adding sufficiently large
positive diagonal blocks, and the minimum rank is attained because feasible
ranks lie in $\{0,\ldots,Lk\}$.  Necessity follows from the compressed input
de Branges--Rovnyak kernel
\[
 X^*\frac{I-S(z)^*S(w)}{1-\overline zw}X.
\]
At the sampled boundary nodes it is positive, has the off-diagonal blocks
\eqref{eq:gp}, and has rank no larger than $\mcdeg S$.

Conversely, factor a feasible rank-$s$ completion as
$P_{ij}=F_i^*F_j$, $F_i\in\C^{s\times k}$.  The defining identity is
equivalent, including on the diagonal, to equality of the Gram matrices of
\begin{equation}\label{eq:lurking}
 \begin{pmatrix}\zeta_iF_i u\\Xu\end{pmatrix}
 \quad\text{and}\quad
 \begin{pmatrix}F_i u\\W_i u\end{pmatrix}.
\end{equation}
Hence their correspondence extends to a unitary colligation
$\left(\begin{smallmatrix}A&B\\C&D\end{smallmatrix}\right)$ on
$\C^s\oplus\C^N$.  Its transfer
\[
 S(z)=D+zC(I-zA)^{-1}B
\]
interpolates the frames.  Any unit-circle eigenspace of the contraction $A$
is reducing for the colligation and decoupled from $B$ and $C$; deleting it
leaves the transfer unchanged and makes it regular at the nodes.  Thus its
degree is at most $s$.  Applying necessity to a minimum-rank completion
precludes a smaller degree and proves \eqref{eq:pick-rank-min}.

\section{Compiler and verification contract}

For a fixed feasible degree $n$, the implementation uses the following
deterministic contract.
\begin{enumerate}
 \item Count positive principal angles and allocate $n-r$ extra subdivisions
 to the principal planes.
 \item Subdivide each $\beta_j$ into positive $\alpha_{j\ell}$.
 \item Initialize $\delta_{j\ell}=2\alpha_{j\ell}$ and distribute the surplus
 $A-2B$ without exceeding $2\pi-4\alpha_{j\ell}$ per gate.
 \item Use \eqref{eq:rhofromdelta} for each scalar zero and
 \eqref{eq:gateoverlap} for each rank-one projector.
 \item Multiply in physical order and evaluate analytic derivatives, never
 finite differences.  Integrate regular cases adaptively and use the exact
 phase antiderivative when a Poisson peak is below numerical resolution.
\end{enumerate}

\begin{table}[h]
\centering
\small
\begin{tabular}{@{}lll@{}}
\toprule
Interface & Independent invariant & Frozen tolerance \\
\midrule
Scalar factor & phase formula vs. adaptive Poisson integral & $2\times10^{-9}$ \\
Rank-one gate & endpoint line angle and unit-vector norm & $5\times10^{-8}$ \\
Potapov product & boundary unitarity and $Q\succeq0$ & $2\times10^{-8}$ \\
Arc compiler & prescribed action and principal angles & $2\times10^{-7}$ \\
Full circle & $\int\tr Q=2\pi n$ & $5\times10^{-7}$ \\
Pick synthesis & colligation unitarity and interpolation & $2\times10^{-7}$ \\
\bottomrule
\end{tabular}
\caption{Numerical contract.  The independent verifier uses separate code
and tighter algebraic checks for the scalar, gate and cycle identities.}
\label{tab:verification}
\end{table}

\bibliographystyle{plain}
\bibliography{references}

\end{document}
